\subsubsection{Control Register (CTRL)}

Address offset: \textbf{0x18}
\newline
Reset value: 0x0000 0000

\begin{figure}[H]
    \centering
    \begin{bytefield}[bitwidth=\bflenhalfword, endianness=big]{16}
        \bitheader[lsb=16]{16-31} \\
        \bitbox{16}{Reserved}
    \end{bytefield}
\end{figure}

\begin{figure}[H]
    \centering
    \begin{bytefield}[bitwidth=\bflenhalfword, endianness=big]{16}
        \bitheader{0-15} \\
        \bitbox{7}{Reserved} &
        \bitbox{1}{WD} &
        \bitbox{7}{Reserved} &
        \bitbox{1}{VEN} \\
        \bitbox[]{7}{} &
        \bitbox[]{1}{R/W} &
        \bitbox[]{7}{} &
        \bitbox[]{1}{R/W}
    \end{bytefield}
\end{figure}

\begin{itemize}
    \item \textbf{Bit 31:9 - Reserved Bits}
    \item \textbf{Bit 8 - WD: Write Data}\\
    Wenn dieses Bit auf 1 gesetzt wird, werden die in VDATR, COLR und ADDRR gesetzten Daten übernommen.
    Dieses Bit wird nach der Übernahme auf 0 zurückgesetzt.
    \item \textbf{Bit 7:1 - Reserved Bits}
    \item \textbf{Bit 0 - VEN: Visualization Enable}\\
    Wenn dieses Bit auf 1 gesetzt ist, werden die in VDATR gesetzten Daten auf der VGA Ausgabe ausgegeben.
\end{itemize}